\documentclass[12pt]{article}

% ==== topological-phases-of-matter : canonical macros (v1, 2026-07-14) ====
% Requires amsmath, amssymb, amsfonts (mathfrak, mathbf, mathcal all standard).
\usepackage{amsmath,amssymb,amsfonts}

% --- condensation / probes ---
\newcommand{\cond}[1]{\underline{#1}}          % condensation X |-> \cond{X}
\newcommand{\cg}[1]{\underline{#1}}            % condensed group (alias, for readability)
\newcommand{\Cont}{\operatorname{Cont}}        % Cont(S,X)
\newcommand{\Shape}{\operatorname{Shape}}      % shape / homotopy type

% --- interaction space & locality ---
\newcommand{\Int}{\mathcal{I}}                 % \Int_{d,G} interaction space
\newcommand{\Ffun}{F}                          % F-function (interaction decay)
\newcommand{\BF}{\mathcal{B}_{F}}              % interaction Banach space for F

% --- moduli stacks ---
\newcommand{\Ham}{\mathfrak{Ham}}              % \Ham_{d,G}
\newcommand{\Gap}{\mathfrak{Gap}}              % \Gap_{d,G}
\newcommand{\Phase}{\mathfrak{Phase}}          % \Phase_{d,G}
\newcommand{\Phases}{\operatorname{Phases}}    % \Phases_{d,G} = pi_0 Shape Phase
\newcommand{\Weq}{\mathcal{W}}                 % class of equivalences to invert
\newcommand{\st}{\mathrm{st}}                  % stabilization superscript
\newcommand{\hfix}[1]{h#1}                     % homotopy-fixed-point superscript, e.g. \Phase_d^{\hfix{\cg{G}}}
\newcommand{\quotstack}[2]{[#1/#2]}            % action stack [Y/\cg{G}]

% --- gap ---
\newcommand{\gapof}[1]{\operatorname{gap}(#1)} % gap(H)
\newcommand{\gapbound}{\Delta}                  % uniform gap bound

% --- stacking / spectrum ---
% \boxtimes is built in (stacking product)
\newcommand{\IPcond}{\mathbf{IP}^{\mathrm{cond}}}  % invertible condensed phase spectrum

% --- transitions / discriminant / charge ---
\newcommand{\Sig}{\Sigma}                       % \Sig_f gapless discriminant
\newcommand{\rc}{\partial\nu}                   % relative transition charge in E^{q+1}(B,U_f)

% --- observables & solid K-theory ---
\newcommand{\Kop}{K_{\mathrm{op}}}             % operator / topological K-theory
\newcommand{\Kalg}{K_{\mathrm{alg}}}           % algebraic K-theory
\newcommand{\Solid}{\operatorname{Solid}}      % solidification functor

% --- disorder ---
\newcommand{\dis}{\Omega}                       % disorder hull \Omega = F^{Z^d}
% ==== end canonical macros ====


% amsthm only (amsmath, amssymb, amsfonts are loaded by the canonical macro block above)
\usepackage{amsthm}

% Diagrams
\usepackage{tikz-cd}
\usepackage{tikz}
\usetikzlibrary{arrows.meta}

% References
\usepackage{thmtools}      % shared-counter theorem envs with correct cleveref names
\usepackage{needspace}     % keep short environments off page breaks
\usepackage{hyperref}
\usepackage{cleveref}

% Graphics
\usepackage{graphicx}

% Page layout
\usepackage[margin=1in]{geometry}

% GrokRxiv sidebar (uses the modern kernel shipout hook, not the obsolete everypage)
\usepackage{xcolor}

% --- paper-local macros (Part II) ---
\newcommand{\qla}{\mathcal{A}}                  % quasi-local C*-algebra
\newcommand{\loc}{\mathrm{loc}}                 % local subalgebra subscript
\newcommand{\St}{\mathcal{S}}                   % state space S(A)
\newcommand{\Pu}{\mathcal{P}}                   % pure states (extreme points)
\newcommand{\Hil}{\mathcal{H}}                  % Hilbert space
\newcommand{\Rep}{\pi}                          % representation
\newcommand{\cyc}{\Omega_{\phi}}                % GNS cyclic vector (context-clear)
\newcommand{\Aut}{\operatorname{Aut}}          % automorphism group
\newcommand{\Prob}{\operatorname{Prob}}        % probability measures
\newcommand{\dalg}{\mathfrak{A}}               % disordered observable algebra
\newcommand{\CAlg}{\mathsf{C^{*}Alg}}          % category of C*-algebras
\newcommand{\CStar}{\mathsf{C^{*}St}}          % C*-algebras with a state
\newcommand{\Cond}{\mathsf{Cond}}              % condensed sets
\newcommand{\RR}{\mathbb{R}}
\newcommand{\CC}{\mathbb{C}}
\newcommand{\ZZ}{\mathbb{Z}}
\newcommand{\NN}{\mathbb{N}}
\newcommand{\TT}{\mathbb{T}}
\newcommand{\Mn}{M_{n}(\CC)}
\newcommand{\wkstar}{\mathrm{w\text{-}\ast}}
\newcommand{\norm}[1]{\lVert #1 \rVert}
\newcommand{\abs}[1]{\lvert #1 \rvert}
\newcommand{\ip}[2]{\langle #1 , #2 \rangle}
\newcommand{\rank}{\operatorname{rank}}
\newcommand{\Tr}{\operatorname{Tr}}
\newcommand{\id}{\mathbf{1}}
\DeclareMathOperator{\ev}{ev}

% --- theorem environments ---
% thmtools 'sibling=theorem' keeps shared numbering (Theorem 2.1, Proposition 2.2, ...)
% while registering the correct cleveref name for each environment, so cross-references
% print "Proposition"/"Lemma"/... rather than mislabelling everything "Theorem".
\declaretheorem[name=Theorem,numberwithin=section]{theorem}
\declaretheorem[name=Proposition,sibling=theorem]{proposition}
\declaretheorem[name=Lemma,sibling=theorem]{lemma}
\declaretheorem[name=Corollary,sibling=theorem]{corollary}
\declaretheorem[name=Definition,style=definition,sibling=theorem]{definition}
\declaretheorem[name=Example,style=definition,sibling=theorem]{example}
\declaretheorem[name=Remark,style=remark,sibling=theorem]{remark}
% program-labelled headline theorems: II-A, II-B, II-C (stable across the series)
\theoremstyle{plain}
\newtheorem{progthm}{Theorem}
\renewcommand{\theprogthm}{II-\Alph{progthm}}
% conjectures with stable II-N numbering (kept fixed across papers)
\newtheorem{conjecture}{Conjecture}
\renewcommand{\theconjecture}{II-\arabic{conjecture}}

% Make hyperref anchors carry the section number: with a section-reset shared
% counter the bare anchors (proposition.2, definition.1, ...) otherwise collide
% across sections and mis-target links.
\makeatletter
\newcommand{\fixtheH}[1]{%
  \@ifundefined{theH#1}{}%
    {\expandafter\renewcommand\csname theH#1\endcsname{\thesection.\arabic{theorem}}}}
\fixtheH{theorem}\fixtheH{proposition}\fixtheH{lemma}\fixtheH{corollary}%
\fixtheH{definition}\fixtheH{example}\fixtheH{remark}
\makeatother

\crefname{progthm}{Theorem}{Theorems}
\Crefname{progthm}{Theorem}{Theorems}
\crefname{conjecture}{Conjecture}{Conjectures}
\Crefname{conjecture}{Conjecture}{Conjectures}
\frenchspacing            % consistent inter-sentence / abbreviation spacing

\hypersetup{
  colorlinks=true,
  linkcolor=blue!55!black,
  citecolor=green!45!black,
  urlcolor=blue!55!black
}

% --- GrokRxiv DOI sidebar ---
\definecolor{grokgray}{RGB}{110,110,110}
\AddToHook{shipout/background}{%
  \ifnum\value{page}=1
    \begin{tikzpicture}[remember picture, overlay]
      \node[
        rotate=90,
        anchor=south,
        font=\Large\sffamily\bfseries\color{grokgray},
        inner sep=0pt
      ] at ([xshift=38pt, yshift=0.52\paperheight]current page.south west)
      {GrokRxiv:2026.07.positivity-cstar-norms\quad
       [\,math.OA\,]\quad
       14 Jul 2026};
    \end{tikzpicture}
  \fi
}

\title{\bfseries Positivity, $C^{*}$-Norms, and Condensed State Spaces of
Quasi-Local Algebras\\[0.4em]
\large Part II of \emph{Topological Phases of Matter in the
Condensed-Mathematics Paradigm}}

\author{Matthew Long \\
\textit{The YonedaAI Collaboration, YonedaAI Research Collective} \\
Chicago, IL \\
\texttt{matthew@yonedaai.com} $\cdot$ \url{https://yonedaai.com}}

\date{14 July 2026}

\begin{document}
\maketitle

\begin{abstract}
The observable side of a quantum lattice system is a quasi-local $C^{*}$-algebra;
its states, its positivity structure, and its $K$-theory are the raw material from
which topological invariants of phases are cut. This paper recasts that material in
condensed mathematics, separating what is provable today from what is not. For a spin
system with finite on-site dimension the
quasi-local algebra $\qla$ is a separable unital approximately finite-dimensional
(AF) algebra, uniformly hyperfinite (UHF) in the homogeneous case, and its
condensation $\cond{\qla}(S)=C(S,\qla)$ is a sheaf of $C^{*}$-algebras on the site
of profinite sets, \emph{light} because $\qla$ is separable. We prove that the
weak-$\ast$ compact convex state space $\St(\qla)$ condenses to a compact Hausdorff
condensed set on which the condensation functor is fully faithful, and that
positivity together with normalization cuts it out as a closed condensed subobject
of the condensed dual ball (Theorem~II-A). The Gelfand--Naimark--Segal construction
is a functor on pointed $C^{*}$-algebras, and on a uniformly gapped family the
ground-state section is weak-$\ast$ continuous by the
Bachmann--Michalakis--Nachtergaele--Sims spectral-flow cocycle, so it is a genuine
point of the condensed state space over the gapped locus, the interface with
Parts~I and~III. The load-bearing bridge is Aoki's solidification theorem:
solidification of the connective algebraic $K$-theory of the condensed algebra
recovers connective semitopological $K$-theory, and, for a real Banach algebra, the
full operator $K$-theory after inverting the Bott class (Theorem~II-B). The
tenfold-way invariants of free-fermion and disordered phases are then values of an
intrinsically condensed-mathematical functor, though the numbers themselves are
unchanged. For disorder we recall Bellissard's
crossed product $C(\dis)\rtimes\ZZ^{d}$ over a profinite hull $\dis=Q^{\ZZ^{d}}$ and
show that its condensation is functorial in the finite quotients of $\dis$
(Theorem~II-C), so profinite probes \emph{match} the physical configuration space.
Five numbered conjectures record what remains open. Accompanying Haskell code
verifies positivity, the GNS construction, and the qubit Bloch ball.
\end{abstract}

\tableofcontents

\section{Introduction}
\label{sec:intro}

\subsection{The problem this paper addresses}

The program that this series develops treats a topological phase of matter as a
connected component of a stabilized condensed stack of uniformly gapped local
Hamiltonians. The prospectus for the program lists five analytic inputs that the
condensed formalism organizes but does not by itself supply: locality and
Lieb--Robinson estimates; positivity and $C^{*}$-norm conditions; existence and
stability of a thermodynamic gap; the passage from lattice models to effective
field theories; and the physical realizability of abstract classes. This is the
paper for the second input.

The word ``positivity'' is doing real work. A quantum lattice system is presented
combinatorially (sites, local Hilbert spaces, a decaying interaction), but its
physical content is carried by an algebra of observables and by the states on that
algebra. That algebra is not an arbitrary Banach algebra: it is a $C^{*}$-algebra,
and the extra axiom $\norm{a^{*}a}=\norm{a}^{2}$ forces a positivity structure that
determines its state space, its representation theory, and ultimately its
$K$-theory. Topological invariants of free-fermion and disordered phases are, in the
end, classes in the operator $K$-theory of this algebra. If the condensed
reformulation is to reach those invariants, it must first carry the $C^{*}$-algebra,
its states, and its positivity into the condensed world without loss. Our task is
to check that this transfer is faithful, to identify what becomes provable once it
is done, and to name the bridge that turns the operator-$K$ invariant into an
intrinsically condensed object.

\subsection{Results and stance}

We hold to the editorial discipline of the series: a statement is labelled
\emph{Theorem} or \emph{Proposition} only when it has a complete proof assembled
from cited, present-day mathematics, and everything else is a numbered
\emph{Conjecture}. Under that discipline the paper proves the following.

\begin{itemize}
\item The quasi-local algebra of a finite-on-site-dimension spin system is a
separable unital AF (UHF in the homogeneous case) $C^{*}$-algebra
(\Cref{prop:uhf}), and its condensation $\cond{\qla}$ is a sheaf of $C^{*}$-algebras
on profinite sets, light because $\qla$ is separable (\Cref{prop:condcstar}).

\item The state space $\St(\qla)$ is weak-$\ast$ compact convex Hausdorff, its
condensation is a compact Hausdorff condensed set on which condensation is fully
faithful, and positivity plus normalization exhibit it as a closed condensed
subobject of the condensed dual ball (\Cref{thm:statespace}, labelled II-A in the
program).

\item GNS is a functor on the category of $C^{*}$-algebras with a state
(\Cref{prop:gns}), and on a uniformly gapped $C^{1}$-family the ground-state
section is weak-$\ast$ continuous (\Cref{thm:groundstate}), hence a point of the
condensed state space over the gapped locus.

\item Aoki's solidification theorem identifies operator $K$-theory of a real Banach
algebra with the solidification of the algebraic $K$-theory of its condensation
(\Cref{thm:aoki}, labelled II-B); consequently the tenfold-way invariants factor
through a condensed-mathematical functor (\Cref{cor:condinvariant}), a statement we
are careful not to inflate.

\item For a profinite disorder hull the covariant observable algebra is Bellissard's
crossed product $C(\dis)\rtimes\ZZ^{d}$, a separable $C^{*}$-algebra whose
$K$-theory carries the disorder-averaged invariants, and its condensation is
functorial in the finite quotients of $\dis$ (\Cref{thm:crossed}, labelled II-C).
\end{itemize}

The open content is organized as \Cref{conj:substack,conj:solidnat,conj:real,conj:classifying,conj:split}:
that positivity and the $C^{*}$-identity cut the Hamiltonian stack out as a closed
condensed substack; that the solid invariant is natural in the disorder hull; that
a Real/$KKO$ refinement recovers the eightfold $KO$-degrees of the Kitaev table;
that the condensed state space is a classifying object for condensed
representations; and that the split property and superselection structure admit
condensed enhancements.

A word on what this paper does \emph{not} claim. It does not produce a new physical
invariant. The Chern number of a Chern insulator, the $\ZZ_{2}$ index of a
time-reversal-invariant insulator, and the winding number of the SSH chain
are unchanged. What changes is their provenance: they are exhibited as values of a
functor whose inputs and intermediate objects are condensed. The value proposition
is organizational, and we say so plainly rather than dress it as a breakthrough.

\subsection{Relation to companion papers}
\label{sec:companions}

This is Part~II of six. Part~I,
\emph{Condensed Locality}~\cite{paperLocality}, builds the interaction Banach space
$\BF$ and shows that the Heisenberg dynamics is a morphism of light condensed sets
$\RR\times\cond{\BF}\to\cond{\Aut(\qla)}$; that action is by automorphisms of the
very algebra whose states we condense here, and the equivalence class $\Weq$ that
Part~I generates through quasi-adiabatic continuation is what makes the ground-state
section of \Cref{thm:groundstate} well defined up to the physically correct notion
of sameness. Part~III,
\emph{The Uniformly Gapped Substack}~\cite{paperGap}, is where the uniform-gap
hypothesis of \Cref{thm:groundstate} is analyzed; our continuous ground-state
section is a section over their gapped locus $U_{f}\subset\Gap_{d,G}$, and the
weak-$\ast$ continuity we prove is the reason $\pi_{0}$ of their localized stack is
well posed on the observable side. Part~IV,
\emph{From Lattice Models to Effective Field Theories}~\cite{paperEFT},
group-completes the invertible sector into the spectrum $\IPcond_{d,G}$; the Aoki
bridge of \Cref{thm:aoki} is precisely the analytic completion their construction
requires, and our \Cref{cor:condinvariant} supplies the observable-algebra
$K$-theory they compare against Freed--Hopkins bordism. Part~V,
\emph{Physical Realizability}~\cite{paperRealizability}, studies which of the
operator-$K$ classes we place in condensed form are realized by gapped lattice
models; our disorder functoriality (\Cref{thm:crossed}, \Cref{conj:solidnat}) is the
input to their functorial-realizability question over disorder hulls. Part~VI,
\emph{A Modular Research Program}~\cite{paperSynthesis}, assembles the modules; this
paper is module~2, the positivity/$C^{*}$ layer, and its emergent contribution to
the synthesis is the functorial solid invariant of disordered families.

\subsection{Notation}

We use the canonical macros of the series: $\cond{X}$ for condensation with
$\cond{X}(S)=\Cont(S,X)$, $\Ham_{d,G}$, $\Gap_{d,G}$, $\Phase_{d,G}$ for the moduli
stacks, $\Weq$ for the inverted equivalences, $\Kop$ and $\Kalg$ for operator and
algebraic $K$-theory, $\Solid$ for solidification, and $\dis=Q^{\ZZ^{d}}$ for the
profinite disorder hull. Paper-local objects are $\qla$ (the quasi-local algebra),
$\St(\qla)$ and $\Pu(\qla)$ (states and pure states), and $\dalg$ (the disordered
observable algebra). Throughout, $C^{*}$-algebras are complex and unital unless we
say otherwise; the sole exception is \Cref{sec:aoki}, where Aoki's theorem is a
statement about \emph{real} algebras and we flag the change explicitly.

\section{Condensed mathematics and quasi-local \texorpdfstring{$C^{*}$}{C*}-algebras}
\label{sec:framework}

We recall only what we use, and we fix conventions so that the later sections read
cleanly. Nothing in this section is new; the point is to make the two categories,
condensed sets and $C^{*}$-algebras, interoperate.

\subsection{Condensed sets and profinite probes}

A \emph{profinite set} is a cofiltered limit of finite sets, equivalently a compact,
Hausdorff, totally disconnected topological space. The category of profinite sets
carries the topology whose covers are the finite jointly surjective families; a
\emph{condensed set} is a sheaf of sets for this topology~\cite{clausen-scholze-condensed,barwick-haine-pyknotic}.
The basic examples come from spaces.

\begin{definition}[Condensation]
\label{def:condensation}
For a topological space $X$, its \emph{condensation} $\cond{X}$ is the condensed set
\[
  \cond{X}(S) \;=\; \Cont(S,X), \qquad S \text{ profinite,}
\]
with restriction along maps of profinite sets given by precomposition.
\end{definition}

The assignment $X\mapsto\cond{X}$ is functorial. On the subcategory of compactly
generated (weak Hausdorff) spaces it is fully faithful, and it preserves finite
products, because $\Cont(S,X\times Y)=\Cont(S,X)\times\Cont(S,Y)$ for every probe
$S$~\cite{clausen-scholze-condensed}. Full faithfulness is the technical reason a
condensed reformulation loses nothing: an ordinary compact space and its
condensation carry the same maps into and out of them.

For the disorder applications we will want a size-controlled variant. A profinite
set is \emph{light} if it is second countable, equivalently a countable cofiltered
limit of finite sets; \emph{light condensed sets} are sheaves on the light profinite
sets~\cite{clausen-scholze-analytic-stacks}. The condensation of a separable
(metrizable, second countable) space is determined by its values on light probes,
and it is convenient to regard it as a light condensed set. Our algebras are
separable, so light condensed structure suffices everywhere below, and we avoid the
set-theoretic subtleties of the unrestricted theory.

\subsection{Quasi-local algebras of spin systems}

Fix a countable set $L$ of sites (for definiteness a lattice $\ZZ^{d}$ with the
graph metric, though only countability and the net of finite subsets matter here).
To each site $x$ attach a finite-dimensional Hilbert space $\Hil_{x}$ with
$\dim\Hil_{x}=n_{x}<\infty$, and set $\Hil_{\Lambda}=\bigotimes_{x\in\Lambda}\Hil_{x}$
for finite $\Lambda\Subset L$. The \emph{local algebra} at $\Lambda$ is the full
matrix algebra
\[
  \qla_{\Lambda} \;=\; \mathcal{B}(\Hil_{\Lambda})
  \;\cong\; \bigotimes_{x\in\Lambda} M_{n_{x}}(\CC).
\]
For $\Lambda\subseteq\Lambda'$ the map $a\mapsto a\otimes\id_{\Lambda'\setminus\Lambda}$
is a unital $\ast$-monomorphism $\qla_{\Lambda}\hookrightarrow\qla_{\Lambda'}$. The
\emph{local subalgebra} is the union
$\qla_{\loc}=\bigcup_{\Lambda\Subset L}\qla_{\Lambda}$, an incomplete normed
$\ast$-algebra, and the \emph{quasi-local algebra} is its completion
\[
  \qla \;=\; \overline{\qla_{\loc}}^{\;\norm{\cdot}},
\]
the $C^{*}$-inductive limit of the net $(\qla_{\Lambda})_{\Lambda}$.

\begin{proposition}[UHF/AF structure of the quasi-local algebra]
\label{prop:uhf}
The quasi-local algebra $\qla$ is a separable unital $C^{*}$-algebra. It is
approximately finite-dimensional: an inductive limit of finite-dimensional
$C^{*}$-algebras with unital connecting maps. If $L$ is infinite and $n_{x}=n$ for
all $x$, then $\qla$ is the uniformly hyperfinite algebra of type $n^{\infty}$,
\[
  \qla \;\cong\; \overline{\bigotimes_{x\in L} M_{n}(\CC)}
        \;=\; M_{n}^{\otimes\infty}.
\]
\end{proposition}

\begin{proof}
Enumerate $L=\{x_{1},x_{2},\dots\}$ and set $\Lambda_{k}=\{x_{1},\dots,x_{k}\}$. The
finite subsets $\Lambda_{k}$ are cofinal in the net of all finite subsets, so
$\qla=\overline{\bigcup_{k}\qla_{\Lambda_{k}}}$ is the closure of an increasing union
of the finite-dimensional $C^{*}$-algebras $\qla_{\Lambda_{k}}$ along the unital
inclusions $\qla_{\Lambda_{k}}\hookrightarrow\qla_{\Lambda_{k+1}}$. By definition
this exhibits $\qla$ as the $C^{*}$-inductive limit
$\varinjlim_{k}\qla_{\Lambda_{k}}$, which is AF. Each $\qla_{\Lambda_{k}}$ is finite
dimensional, hence separable, and it contains the countable
$\mathbb{Q}(i)$-$\ast$-subalgebra of matrices with Gaussian-rational entries; the
union of these over $k$ is a countable norm-dense $\ast$-subalgebra of $\qla$, so
$\qla$ is separable. The unit of $\qla_{\Lambda_{1}}$ is a unit for the limit. When
$n_{x}=n$ for all $x$, each $\qla_{\Lambda_{k}}\cong M_{n}(\CC)^{\otimes k}\cong
M_{n^{k}}(\CC)$ and the connecting maps are the standard unital embeddings
$M_{n^{k}}\hookrightarrow M_{n^{k+1}}$, $a\mapsto a\otimes\id_{n}$; the inductive
limit of this system is by definition the UHF algebra of type $n^{\infty}$.
\end{proof}

The operator-algebraic viewpoint on gapped phases takes exactly this $\qla$ as the
algebra of quasi-local observables; the ground states of physical interest are
states on $\qla$~\cite{ogata-classification-review,nsy-bulkgap}. That $\qla$ is AF,
and in particular nuclear and separable, is what will make its condensation light
and its state space metrizable.

\subsection{Condensation of a \texorpdfstring{$C^{*}$}{C*}-algebra}
\label{sec:condofalg}

A $C^{*}$-algebra is an algebra object with involution in the symmetric monoidal
category of (compactly generated) topological spaces: multiplication, addition,
scalar action, and involution are continuous, and the norm makes the underlying set
a complete metric space. Since condensation preserves finite products, it carries
this structure across.

\begin{proposition}[Condensation is a sheaf of $C^{*}$-algebras; separable $\Rightarrow$ light]
\label{prop:condcstar}
Let $A$ be a $C^{*}$-algebra. The condensation $\cond{A}$, $\cond{A}(S)=C(S,A)$, is a
$C^{*}$-algebra object in condensed sets: for every profinite $S$ the value $C(S,A)$
is a $C^{*}$-algebra under pointwise operations and the supremum norm, and the
restriction maps are $\ast$-homomorphisms. If $A$ is unital, so is $\cond{A}$. If $A$
is separable, $\cond{A}$ is a light condensed set.
\end{proposition}

\begin{proof}
That $S\mapsto\Cont(S,A)$ is a sheaf, i.e.\ a condensed set, is the defining
property of the condensation of the topological space
$A$~\cite{clausen-scholze-condensed}. The $C^{*}$-operations of $A$ are continuous
maps $A\times A\to A$ and $A\to A$; post-composition sends a continuous $S$-family
to a continuous $S$-family, and because condensation preserves products these
assemble into morphisms of condensed sets
$\cond{A}\times\cond{A}\to\cond{A}$ and $\cond{A}\to\cond{A}$ satisfying the algebra
and involution axioms levelwise. Concretely, for compact $S$ the space $C(S,A)$ with
$(f\cdot g)(s)=f(s)g(s)$, $f^{*}(s)=f(s)^{*}$, and $\norm{f}=\sup_{s}\norm{f(s)}$ is
a $C^{*}$-algebra: it is a Banach $\ast$-algebra, and the $C^{*}$-identity is
inherited pointwise,
\[
  \norm{f^{*}f} = \sup_{s}\norm{f(s)^{*}f(s)}
               = \sup_{s}\norm{f(s)}^{2}
               = \Big(\sup_{s}\norm{f(s)}\Big)^{2} = \norm{f}^{2}.
\]
Restriction along $\theta:S'\to S$ is $f\mapsto f\circ\theta$, evidently a unital
$\ast$-homomorphism. When $A$ is unital the constant function at $\id_{A}$ is a unit
for each $C(S,A)$, compatible with restriction. Finally, if $A$ is separable it is
second countable, so its condensation is determined by its restriction to light
profinite probes and is a light condensed
set~\cite{clausen-scholze-analytic-stacks}.
\end{proof}

\begin{remark}
The sheaf condition here is descent for $A$-valued continuous functions along finite
jointly surjective families $\{S_{i}\to S\}$: such a family gives a topological
quotient $\coprod_{i}S_{i}\twoheadrightarrow S$ of compact Hausdorff spaces, and a
continuous $A$-valued function on $S$ is the same as a compatible family on the
$S_{i}$ agreeing on the fibered products $S_{i}\times_{S}S_{j}$. This is the reason
condensation, and not merely the underlying set, is the right target: it remembers
how continuous families glue.
\end{remark}

\section{Positivity and the condensed state space}
\label{sec:states}

We now condense the state space and read off positivity as a closed condition. This
is the first place where the $C^{*}$-axiom, and not just the Banach structure, is
used.

\subsection{States, positivity, and weak-\texorpdfstring{$\ast$}{*} compactness}

Let $A$ be a unital $C^{*}$-algebra. A \emph{state} is a linear functional
$\phi:A\to\CC$ that is positive, $\phi(a^{*}a)\ge 0$ for all $a$, and normalized,
$\phi(\id)=1$. Positivity forces continuity with $\norm{\phi}=\phi(\id)=1$, so states
lie in the unit ball $(A^{*})_{1}$ of the dual. Write $\St(A)$ for the set of states
and $\Pu(A)$ for its extreme points, the \emph{pure states}.

\begin{lemma}[Weak-$\ast$ compactness]
\label{lem:alaoglu}
$\St(A)$ is a weak-$\ast$ closed, convex subset of $(A^{*})_{1}$, hence weak-$\ast$
compact and Hausdorff. If $A$ is separable, $\St(A)$ is metrizable, so it is a
compact convex Polish space, and $\St(A)=\overline{\operatorname{conv}}\,\Pu(A)$.
\end{lemma}

\begin{proof}
By Banach--Alaoglu $(A^{*})_{1}$ is weak-$\ast$ compact and Hausdorff. For each fixed
$a\in A$ the evaluation $\ev_{a}:\phi\mapsto\phi(a)$ is weak-$\ast$ continuous, so
the sets $\{\phi:\phi(a^{*}a)\ge 0\}=\ev_{a^{*}a}^{-1}([0,\infty))$ and
$\{\phi:\phi(\id)=1\}=\ev_{\id}^{-1}(\{1\})$ are weak-$\ast$ closed. Their
intersection over all $a$ is $\St(A)$, which is therefore weak-$\ast$ closed in the
compact set $(A^{*})_{1}$, hence compact; convexity is immediate from linearity of
the conditions. When $A$ is separable the weak-$\ast$ topology on $(A^{*})_{1}$ is
metrizable, so the closed subset $\St(A)$ is a compact metric space; it is convex
and compact, and Krein--Milman gives
$\St(A)=\overline{\operatorname{conv}}\,\Pu(A)$.
\end{proof}

\begin{example}[Matrix algebras and the Bloch ball]
\label{ex:bloch}
For $A=\Mn$ every state is $\phi_{\rho}(a)=\Tr(\rho a)$ for a unique density matrix
$\rho\ge 0$ with $\Tr\rho=1$, so $\St(M_{n})$ is the set of density matrices and
$\Pu(M_{n})$ is the set of rank-one projections, a copy of $\CC P^{n-1}$. For $n=2$,
writing $\rho=\tfrac12(\id+\vec r\cdot\vec\sigma)$ with $\vec\sigma$ the Pauli
matrices, positivity is exactly $\abs{\vec r}\le 1$: the state space $\St(M_{2})$ is
the closed unit ball $B^{3}\subset\RR^{3}$ (the Bloch ball) and $\Pu(M_{2})$ is its
boundary sphere $S^{2}=\CC P^{1}$. The accompanying Haskell code takes exactly this
picture as its concrete testbed.
\end{example}

\subsection{The condensed state space}

\begin{progthm}[Condensed state space]
\label{thm:statespace}
Let $A$ be a unital $C^{*}$-algebra. Then $\cond{\St(A)}$ is a compact Hausdorff
condensed set, and the condensation functor is fully faithful on it. Positivity and
normalization exhibit $\cond{\St(A)}$ as a closed condensed subobject of the
condensed dual ball $\cond{(A^{*})_{1}}$: for every profinite $S$,
\[
  \cond{\St(A)}(S)
  = \bigl\{\, f\in C\!\bigl(S,(A^{*})_{1}\bigr) :
      f(s)(a^{*}a)\ge 0 \ \text{and}\ f(s)(\id)=1
      \ \text{for all } s\in S,\ a\in A \,\bigr\},
\]
a weak-$\ast$ closed subset of $C(S,(A^{*})_{1})$. If $A$ is separable the subobject
is a closed inclusion of light condensed sets.
\end{progthm}

\begin{proof}
By \Cref{lem:alaoglu}, $\St(A)$ is a compact Hausdorff space, so its condensation is
a compact Hausdorff condensed set, and full faithfulness of $\cond{(-)}$ on compact
Hausdorff spaces is~\cite{clausen-scholze-condensed}. For a profinite $S$ a point of
$\cond{\St(A)}(S)$ is a continuous map $f:S\to\St(A)$, i.e.\ a continuous map
$f:S\to(A^{*})_{1}$ each of whose values is a state; this is the displayed set. Each
defining condition is closed: for fixed $a$ and $s$ the map $f\mapsto f(s)(a^{*}a)$
is continuous on $C(S,(A^{*})_{1})$ (composite of evaluation at $s$ and $\ev_{a^{*}a}$),
so $\{f:f(s)(a^{*}a)\ge 0\}$ is closed, and likewise for normalization; the
intersection over $s\in S$ and $a\in A$ is closed. Because the inclusion
$\St(A)\hookrightarrow(A^{*})_{1}$ is a closed embedding of compact Hausdorff spaces
and condensation preserves closed embeddings of such spaces, the induced map
$\cond{\St(A)}\to\cond{(A^{*})_{1}}$ is a closed inclusion of condensed sets. When
$A$ is separable both spaces are second countable, so both condensations are light
by \Cref{prop:condcstar} and the inclusion is one of light condensed sets.
\end{proof}

\begin{remark}[Why closedness is the right statement]
A condensed reformulation of ``the state space'' has to answer: what is a
\emph{family} of states over a parameter, and when does a limit of states stay a
state? \Cref{thm:statespace} answers both at once. A family over $S$ is a point of
$\cond{\St(A)}(S)$; positivity is preserved under the profinite probes because it is
a closed condition, so a profinite limit of states is a state. This is the condensed
avatar of the elementary fact that the state space is weak-$\ast$ closed, upgraded
to a statement about all profinite families simultaneously.
\end{remark}

\subsection{Positivity as a condensed order structure}
\label{sec:cone}

It is worth saying what ``positivity is a closed condition'' means at the level of
the algebra, not just the dual, because that is where the $C^{*}$-identity enters and
because it is the form the condition takes when one tries to characterize
$\Ham_{d,G}$ intrinsically (\Cref{conj:substack}). The positive cone of a
$C^{*}$-algebra $A$ is
\[
  A_{+} \;=\; \{\, a^{*}a : a\in A \,\} \;=\; \{\, a\in A_{\mathrm{sa}} : \sigma(a)\subseteq[0,\infty) \,\},
\]
a closed convex cone in the real Banach space $A_{\mathrm{sa}}$ of self-adjoint
elements, with $A_{+}\cap(-A_{+})=\{0\}$; together with the order unit $\id$ it makes
$(A_{\mathrm{sa}},A_{+},\id)$ an order-unit space whose order-unit state space is
exactly $\St(A)$.

\begin{proposition}[Condensation preserves the positive cone and order unit]
\label{prop:cone}
Let $A$ be a unital $C^{*}$-algebra. The condensation $\cond{A_{+}}$ is a closed
sub-condensed-set of $\cond{A_{\mathrm{sa}}}$, and for every profinite $S$,
\[
  \cond{A_{+}}(S) \;=\; \{\, f\in C(S,A_{\mathrm{sa}}) : f(s)\in A_{+} \text{ for all } s\in S \,\}
  \;=\; C(S,A_{+}).
\]
The order unit condenses to the constant section $\id$, and a state on $A$ is the
same as a morphism of condensed sets $\cond{A_{\mathrm{sa}}}\to\cond{\RR}$ that is
$\RR$-linear on probes, carries $\cond{A_{+}}$ into $\cond{\RR_{\ge 0}}$, and sends
$\id$ to $1$.
\end{proposition}

\begin{proof}
The cone $A_{+}$ is a closed subset of $A_{\mathrm{sa}}$, so a continuous map
$f:S\to A_{\mathrm{sa}}$ takes values in $A_{+}$ if and only if $f(s)\in A_{+}$ for
every $s$; this is the displayed equality $\cond{A_{+}}(S)=C(S,A_{+})$, and
closedness of the inclusion into $\cond{A_{\mathrm{sa}}}$ follows as in
\Cref{thm:statespace}. That condensation carries the order unit to the constant
section is immediate. For the last clause, a probe-wise $\RR$-linear morphism
$\cond{A_{\mathrm{sa}}}\to\cond{\RR}$ is, by full faithfulness of condensation on the
compactly generated space $A_{\mathrm{sa}}$, a continuous $\RR$-linear functional
$\phi:A_{\mathrm{sa}}\to\RR$; the conditions $\phi(A_{+})\subseteq\RR_{\ge0}$ and
$\phi(\id)=1$ are precisely those making $\phi$ extend to a state on $A$.
\end{proof}

This is the internal counterpart of \Cref{thm:statespace}: positivity lives on the
algebra as a closed condensed cone, and a state is a cone- and unit-preserving
condensed functional. The $C^{*}$-identity is what pins this cone to the spectrum,
and it is the relation \Cref{conj:substack} would use to cut $\Ham_{d,G}$ out of a
formal-interaction stack.

\section{The GNS construction as a functor}
\label{sec:gns}

The link between a state and a Hilbert-space representation is the
Gelfand--Naimark--Segal construction. We record its functoriality (elementary, but
what lets it interact with the condensed structure) and then
prove the continuity statement that matters for phases: on a uniformly gapped
family, the ground-state section into $\St(\qla)$ is weak-$\ast$ continuous.

\subsection{Functoriality of GNS}

\begin{definition}[The category of pointed $C^{*}$-algebras]
Let $\CStar$ be the category whose objects are pairs $(A,\phi)$ of a unital
$C^{*}$-algebra and a state, and whose morphisms $(A,\phi)\to(B,\psi)$ are unital
$\ast$-homomorphisms $\theta:A\to B$ with $\psi\circ\theta=\phi$ (state-preserving
maps).
\end{definition}

Given $(A,\phi)$, the sesquilinear form $\ip{a}{b}_{\phi}=\phi(a^{*}b)$ is positive
semidefinite; its kernel $N_{\phi}=\{a:\phi(a^{*}a)=0\}$ is a closed left ideal (by
the Cauchy--Schwarz inequality $\abs{\phi(a^{*}b)}^{2}\le\phi(a^{*}a)\phi(b^{*}b)$);
the completion $\Hil_{\phi}=\overline{A/N_{\phi}}$ is a Hilbert space; left
multiplication descends to a $\ast$-representation $\Rep_{\phi}(a)[b]=[ab]$; and the
class $\cyc=[\id]$ is a cyclic unit vector with
$\phi(a)=\ip{\cyc}{\Rep_{\phi}(a)\cyc}$.

\begin{proposition}[GNS is a functor]
\label{prop:gns}
The assignment $(A,\phi)\mapsto(\Hil_{\phi},\Rep_{\phi},\cyc)$ extends to a functor
from $\CStar$ to the category of pointed cyclic representations with intertwining
isometries. A state-preserving map $\theta:(A,\phi)\to(B,\psi)$ induces the isometry
\[
  V_{\theta}:\Hil_{\phi}\to\Hil_{\psi}, \qquad V_{\theta}[a]_{\phi}=[\theta(a)]_{\psi},
\]
satisfying $V_{\theta}\cyc=\Omega_{\psi}$ and
$V_{\theta}\,\Rep_{\phi}(a)=\Rep_{\psi}(\theta(a))\,V_{\theta}$, with
$V_{\theta'\circ\theta}=V_{\theta'}V_{\theta}$ and $V_{\mathrm{id}}=\mathrm{id}$.
\end{proposition}

\begin{proof}
Because $\psi\circ\theta=\phi$ we have
$\ip{\theta(a)}{\theta(b)}_{\psi}=\psi(\theta(a)^{*}\theta(b))
=\psi(\theta(a^{*}b))=\phi(a^{*}b)=\ip{a}{b}_{\phi}$, so $\theta$ maps $N_{\phi}$
into $N_{\psi}$ and the map $[a]_{\phi}\mapsto[\theta(a)]_{\psi}$ is a well-defined
isometry on the dense subspaces, extending to an isometry $V_{\theta}$ of the
completions. Evaluating on the unit gives $V_{\theta}\Omega_{\phi}=\Omega_{\psi}$,
and for $a,b\in A$,
$V_{\theta}\Rep_{\phi}(a)[b]_{\phi}=V_{\theta}[ab]_{\phi}=[\theta(a)\theta(b)]_{\psi}
=\Rep_{\psi}(\theta(a))[\theta(b)]_{\psi}=\Rep_{\psi}(\theta(a))V_{\theta}[b]_{\phi}$,
i.e.\ $V_{\theta}$ intertwines. Functoriality in $\theta$ is immediate from the
definition on representatives.
\end{proof}

\subsection{Continuity of the ground-state section on gapped families}

The physically decisive continuity statement is not about GNS of a single state but
about a \emph{family} of ground states. We use the automorphic-equivalence theorem
for gapped phases, which is the interface with Parts~I and~III.

\begin{theorem}[Weak-$\ast$ continuity of the ground-state section]
\label{thm:groundstate}
Let $s\mapsto\Phi_{s}$, $s\in[0,1]$, be a family of interactions that is $C^{1}$ in
the interaction Banach norm $\BF$ of Part~I, and suppose the family is
\emph{uniformly gapped}: the thermodynamic-limit dynamics has a unique ground state
$\omega_{s}$ on $\qla$ separated from the rest of the spectrum by a gap
$\ge\gapbound>0$ uniformly in $s$. Then there is a strongly continuous family of
quasi-local automorphisms $\alpha_{s}\in\Aut(\qla)$ with
$\omega_{s}=\omega_{0}\circ\alpha_{s}$, and for every local observable
$a\in\qla_{\loc}$ the map $s\mapsto\omega_{s}(a)$ is continuous. Consequently
$s\mapsto\omega_{s}$ is a weak-$\ast$ continuous map $[0,1]\to\St(\qla)$; taking a
profinite probe $S\to[0,1]$ landing in the gapped locus, the ground state is a point
of $\cond{\St(\qla)}(S)$.
\end{theorem}

\begin{proof}
Under the stated hypotheses the Bachmann--Michalakis--Nachtergaele--Sims theorem on
automorphic equivalence within a gapped phase produces the quasi-adiabatic spectral-flow
cocycle: a strongly continuous family of $\ast$-automorphisms $\alpha_{s}$ of $\qla$,
generated by a time-dependent quasi-local interaction built from $\dot\Phi_{s}$ and
the gap, with $\omega_{s}=\omega_{0}\circ\alpha_{s}$~\cite{bmns-automorphic,hastings-wen-qac}.
The generator is quasi-local with Lieb--Robinson-controlled tails, so for a local
observable $a$ the orbit $s\mapsto\alpha_{s}(a)$ is norm-continuous (indeed
$C^{1}$)~\cite{nsy-quasilocality-1,paperLocality}. Hence
$s\mapsto\omega_{s}(a)=\omega_{0}(\alpha_{s}(a))$ is continuous, being the
composition of a norm-continuous $\qla$-valued map with the bounded functional
$\omega_{0}$. Local observables are norm-dense in $\qla$ and each $\omega_{s}$ is a
state with $\norm{\omega_{s}}=1$; a uniformly bounded net of functionals that
converges on a dense set converges weak-$\ast$, so continuity on $\qla_{\loc}$
upgrades to weak-$\ast$ continuity on all of $\qla$. Thus $s\mapsto\omega_{s}$ is a
continuous map into $\St(\qla)$ with its weak-$\ast$ topology. For a profinite
$S$ with a continuous map $S\to[0,1]$ whose image lies in the gapped locus, the
composite $S\to\St(\qla)$ is continuous, i.e.\ an element of $\cond{\St(\qla)}(S)$ by
\Cref{thm:statespace}.
\end{proof}

\begin{remark}[What is used, and what is deferred]
The theorem uses the uniform gap as a hypothesis; whether a given family satisfies it
is the subject of Part~III, and in general the question is undecidable, so we do not
attempt a criterion here. The output is exactly the datum the program needs on the
observable side: a continuous section of states over the gapped locus, which is what
makes the phase label $\nu_{f}$ locally constant and the whole $\pi_{0}$ story well
posed. The representation-theoretic refinement (that GNS assembles into a
continuous field of Hilbert spaces over the base) is delicate because the GNS
representation can jump in dimension, and we record the clean version we can prove
(continuity of the vector-state section) rather than overclaim a continuous field.
The full statement is part of \Cref{conj:classifying}.
\end{remark}

\section{Free fermions and the real observable algebra}
\label{sec:freefermion}

Before invoking Aoki's theorem we make concrete the algebra to which it will be
applied. Free-fermion systems are the best-developed setting for topological
classification, and they are where the \emph{real} structure that Aoki's theorem
requires is visible on the nose. This section fixes the observable algebra, its
quasi-free ground states, and the real structure imposed by the antiunitary
symmetries, and it identifies the SSH winding number as an operator-$K$ class so that
the corollary of \Cref{sec:aoki} lands on a familiar object.

\subsection{The CAR algebra and quasi-free ground states}

For a free-fermion system with single-particle Hilbert space $\mathfrak{h}=\ell^{2}(L)$
(one orbital per site, for simplicity) the observable algebra is the algebra of
canonical anticommutation relations $\mathrm{CAR}(\mathfrak{h})$, the unital
$C^{*}$-algebra generated by $\{a(v):v\in\mathfrak{h}\}$ with $v\mapsto a(v)^{*}$
linear and
\[
  \{a(v),a(w)^{*}\}=\ip{v}{w}\id, \qquad \{a(v),a(w)\}=0 ,
\]
where the inner product is linear in its second argument (the physics convention),
so that $v\mapsto a(v)$ is antilinear and $v\mapsto a(v)^{*}$ is linear.
For separable infinite-dimensional $\mathfrak{h}$ the algebra
$\mathrm{CAR}(\mathfrak{h})$ is the UHF algebra of type $2^{\infty}$; in particular it
is a special case of the quasi-local algebra of \Cref{prop:uhf}, and everything in
\Cref{sec:framework,sec:states,sec:gns} applies to it verbatim.

A quadratic Hamiltonian is specified by a self-adjoint single-particle operator
$h=h^{*}$ on $\mathfrak{h}$, and when $h$ has a spectral gap at the Fermi level $0$ the
ground state is the \emph{quasi-free} state $\omega_{P}$ determined by the Fermi
projection $P=\chi_{(-\infty,0)}(h)$, through the two-point function
\[
  \omega_{P}\bigl(a(v)^{*}a(w)\bigr)=\ip{w}{Pv},
\]
with Wick's theorem fixing all higher correlators. Two facts connect this to the earlier
sections. First, $\omega_{P}$ is a state on $\mathrm{CAR}(\mathfrak{h})$, hence a point
of $\St(\mathrm{CAR}(\mathfrak{h}))$, and the map $P\mapsto\omega_{P}$ is weak-$\ast$
continuous in $P$; a uniformly gapped family of single-particle Hamiltonians thus
gives a continuous ground-state section exactly as in \Cref{thm:groundstate}, and one
does not even need the full automorphic-equivalence machinery in the quasi-free case
because the section is visibly continuous in the Fermi projection. Second, the spectral
gap forces $P$ to have exponentially decaying off-diagonal, so $P$ is a
\emph{controlled} (quasi-local) projection --- the free-fermion shadow of the
clustering established from a gap in Part~I and used throughout Part~III.

\subsection{Real structure and the tenfold way}

The tenfold way arranges gapped free-fermion phases by their behavior under the
antiunitary symmetries: time reversal $T$, particle--hole $C$, and their product, the
chiral symmetry $S=TC$. Implemented on $\mathfrak{h}$, these are antilinear (for $T,C$)
or linear (for $S$) operators squaring to $\pm\id$, and they endow the observable
algebra with a \emph{Real} structure in the sense of Atiyah: a real form
$A^{\RR}\subseteq\mathrm{CAR}(\mathfrak{h})$ (or a $\ZZ_{2}$-graded, Clifford-module
refinement thereof) whose operator $K$-theory, in the degree fixed by the symmetry
class, is the group of the Kitaev periodic
table~\cite{kitaev-periodic-table,thiang-ktheory,prodan-sb}. The eight real classes are
carried by $KO$, the two complex classes by $KU$, and the eightfold Bott periodicity of
$KO$ is the eightfold periodicity of the table. The point for us is structural: the
algebra that classifies free-fermion phases is a \emph{real} $C^{*}$-algebra, and its
invariant is an operator-$K$ class, precisely the input that Aoki's theorem consumes.

\begin{example}[The SSH chain as a $K_{1}$-class]
\label{ex:ssh}
The Su--Schrieffer--Heeger chain~\cite{ssh-1979} is the two-band model with
single-particle Bloch Hamiltonian
\[
  H(k)=\bigl(t_{1}+t_{2}\cos k\bigr)\sigma_{x}+t_{2}\sin k\,\sigma_{y},
  \qquad q(k)=t_{1}+t_{2}e^{ik},
\]
chiral-symmetric under $\sigma_{z}H(k)\sigma_{z}=-H(k)$ and gapped precisely when
$\abs{t_{1}}\neq\abs{t_{2}}$. Chiral symmetry flattens the Fermi projection into the
off-diagonal unitary $u(k)=q(k)/\abs{q(k)}$, and the winding number
\[
  \nu=\frac{1}{2\pi i}\oint u(k)^{-1}\,du(k)\in\ZZ
\]
is the homotopy class of $u:S^{1}\to U(1)$. In operator-algebraic terms $\nu$ is the
class $[u]\in K_{1}(A)$ of the observable algebra $A$ --- the translation-invariant
algebra $C(S^{1})$ in the clean case, or the real-space (possibly disordered) chiral
algebra in general --- and it is computed as the index of the associated Toeplitz
operator, a genuine operator-$K$ invariant~\cite{prodan-sb,bourne-kellendonk-rennie}.
It takes $\nu=1$ for $\abs{t_{1}}<\abs{t_{2}}$ and $\nu=0$ for $\abs{t_{1}}>\abs{t_{2}}$.
This is the flagship example of the program's SSH transition, and it will be the
example on which the corollary of \Cref{sec:aoki} is read.
\end{example}

The SSH class sits in $K_{1}$; the two-dimensional Chern insulator's Hall conductance
sits in $K_{0}$ and is an integer multiple of $e^{2}/h$~\cite{tknn-1982}; the
disordered analogues sit in the $K$-theory of the crossed product of \Cref{sec:disorder}.
In every case the invariant is a class in the operator $K$-theory of a real (or Real)
observable algebra, and that is all \Cref{thm:aoki} needs.

\section{Solidification: operator \texorpdfstring{$K$}{K}-theory as a condensed invariant}
\label{sec:aoki}

We come to the load-bearing result. Topological invariants of free-fermion and
disordered phases are classes in the operator $K$-theory of a Banach or
$C^{*}$-algebra of observables. Aoki's theorem shows that this operator $K$-theory is
recovered from an intrinsically condensed operation applied to the \emph{algebraic}
$K$-theory of the condensed algebra. This section states that theorem with its
hypotheses and draws the consequence for phases, taking care not to inflate it.

\subsection{Aoki's theorem}

We change convention for this section only: algebras are \emph{real}. This is not a
technicality to be waved away. It is exactly the setting in which the physically
relevant $KO$-graded invariants live, so the real hypothesis is a feature. For a
real Banach algebra $A$ its condensation $\cond{A}$, $\cond{A}(S)=C(S,A)$, is a
condensed (indeed solid) ring in the sense of Clausen--Scholze, and one may form its
\emph{connective} algebraic $K$-theory $\Kalg(\cond{A})$ as a condensed spectrum. The
solidification functor $\Solid$ is the reflection onto solid condensed spectra
of~\cite{clausen-scholze-condensed,clausen-scholze-analytic}.

\begin{progthm}[Aoki's solidification theorem~{\cite{aoki-solidification}}]
\label{thm:aoki}
Let $A$ be a real associative algebra and $\cond{A}$ its condensation. The
solidification of the \emph{connective} algebraic $K$-theory of $\cond{A}$ is discrete
and recovers the connective semitopological $K$-theory of $A$ in the sense of
Friedlander--Walker (and Blanc),
\[
  \Solid\bigl(\Kalg(\cond{A})\bigr) \;\simeq\; K^{\mathrm{semi}}(A).
\]
If moreover $A$ is a real Banach algebra, this is the \emph{connective part} of the
topological (operator) $K$-theory; the full, Bott-periodic operator $K$-theory is
obtained only after inverting the Bott class $\beta$,
\[
  \Kop(A) \;\simeq\; \Solid\bigl(\Kalg(\cond{A})\bigr)\bigl[\beta^{-1}\bigr],
\]
an equivalence of spectra, where $\beta$ is the generator of the relevant Bott group
($\beta\in K_{2}^{\mathrm{top}}(\CC)\cong\ZZ$ in the complex case; the order-eight
real Bott class in the $KO$ case).
\end{progthm}

Two points of bookkeeping matter, and we make them explicit rather than bury them.
First, the solidification is of the \emph{connective} algebraic $K$-theory, so it
lands in a connective spectrum with no homotopy in negative degrees, whereas operator
$K$-theory is Bott-periodic and does have negative-degree groups. The two therefore
agree only in nonnegative degrees; the equivalence with the full $\Kop(A)$ holds after
the Bott inversion $[\beta^{-1}]$, and writing
$\Kop(A)\simeq\Solid(\Kalg(\cond{A}))$ without it would be false in negative degrees.
Second, the real versus Real ($KR$) distinction: $\beta$ is the order-eight real Bott
class, and the eightfold periodicity of $KO$ is the eightfold structure of the real
symmetry classes. A word on the topology in the first clause: $A$ carries its given
topology --- discrete when none is specified, in which case $\cond{A}$ is the
condensation of the underlying discrete algebra, whose value on a profinite $S$ is the
algebra of locally constant $A$-valued functions, a nontrivial condensed ring because
profinite sets are richly disconnected, so the semitopological $K$-theory it computes
is nontrivial even in the discrete case. The Banach clause, with its Bott inversion,
is the one we use. The proof is Aoki's and we claim no part of it; what we contribute
is the reading of the theorem inside the program.

\subsection{Free-fermion invariants are condensed-mathematical}

\begin{corollary}[Provenance of the tenfold-way invariants]
\label{cor:condinvariant}
Let a free-fermion or disordered topological insulator or superconductor have real
observable Banach algebra $A$ --- the real form of the (bulk or disorder-averaged)
observable algebra in its symmetry class. Then its tenfold-way operator-$K$
invariant $\nu\in\Kop(A)$ is the value at $A$ of the composite functor
\[
  A \;\longmapsto\; \cond{A} \;\longmapsto\; \Kalg(\cond{A})
    \;\longmapsto\; \Solid\bigl(\Kalg(\cond{A})\bigr)\bigl[\beta^{-1}\bigr]
    \;\simeq\; \Kop(A) \ni \nu.
\]
In particular the invariant factors through condensation, solidification, and the
Bott inversion.
\end{corollary}

\noindent Concretely, for the SSH chain of \Cref{ex:ssh} the winding number
$\nu=[u]\in K_{1}(A)$ is a class in $\Solid(\Kalg(\cond{A}))[\beta^{-1}]$: the integer
that distinguishes the two SSH phases is the value of the condensed-mathematical
composite above. The number is the same integer computed in 1979; what
\Cref{thm:aoki} adds is that it is the output of solidification, followed by Bott
inversion, applied to the algebraic $K$-theory of a condensed algebra.

\begin{proof}
By the $K$-theoretic classification of free-fermion and disordered phases, the
symmetry-class invariant of such a system is a class in the operator $K$-theory of
its (real/Real) observable algebra $A$, in the appropriate $KO$- or
$KU$-degree~\cite{kitaev-periodic-table,thiang-ktheory,prodan-sb}. Apply
\Cref{thm:aoki} to $A$: the operator $K$-theory group in which $\nu$ lives is
$\Solid(\Kalg(\cond{A}))[\beta^{-1}]$ (the Bott-inverted, periodic value), so $\nu$ is
a class in the target of the displayed composite, and the composite computes it.
\end{proof}

\begin{remark}[Exactly how much this claims]
It claims that the invariant's \emph{home} and \emph{construction route} are
condensed-mathematical, not that the invariant is new or numerically different. A
Chern number computed as a solid $K$-class equals the Chern number computed by
integrating a curvature form. The content is that the two mature theories (the
operator $K$-theory of topological insulators and the condensed/solid $K$-theory of
Clausen--Scholze and Aoki) meet here, so the family-theoretic, descent-theoretic,
and profinite-disorder machinery of the condensed side becomes available for an
invariant that was previously computed by hand. This is the sense in which the
program's value is organizational. We state it this way in deference to the honest
appraisal that the program does not, on its own, produce a different Chern number.
\end{remark}

\begin{remark}[Why the real setting is the right one]
The Kitaev periodic table is $KO$-graded: the eightfold Bott periodicity of real
$K$-theory is the eightfold structure of the real symmetry classes, with the two
complex classes carried by $KU$. Aoki's theorem is proved for real algebras, so it
lands precisely where the physics lives; a complex-only statement would miss the
time-reversal and particle--hole classes. The refinement to a genuinely Real ($KR$)
statement matching all ten classes is \Cref{conj:real}.
\end{remark}

\section{Disorder, crossed products, and profinite probes}
\label{sec:disorder}

Disorder is where the condensed language stops repackaging manifolds and starts
matching the physical configuration space. A disordered family is literally a point
of the Hamiltonian moduli object over a profinite base.

\subsection{The disorder hull as a profinite probe}

Let $Q$ be a finite set of local configurations\footnote{The prospectus and the program's notation write this alphabet as $F$; we use $Q$ to avoid collision with the Nachtergaele--Sims--Young $F$-function of Part~I (carried here only inside the interaction Banach space $\BF$).} and
\[
  \dis \;=\; Q^{\ZZ^{d}},
\]
the space of configurations, with the product topology. It is compact, Hausdorff,
and totally disconnected (a profinite set), and the shift action of $\ZZ^{d}$ by
translations makes it a profinite $\ZZ^{d}$-space. Restricting a configuration to a
finite window $\Lambda_{i}\Subset\ZZ^{d}$ gives finite quotients
$\dis\twoheadrightarrow\dis_{i}=Q^{\Lambda_{i}}$, and $\dis=\varprojlim_{i}\dis_{i}$.

\begin{proposition}[$C(\dis)$ is a commutative AF algebra]
\label{prop:hullAF}
For $\dis=\varprojlim_{i}\dis_{i}$ profinite, the pullbacks along
$\dis\to\dis_{i}$ induce an isometric $\ast$-isomorphism
\[
  C(\dis) \;\cong\; \varinjlim_{i} C(\dis_{i}),
\]
so $C(\dis)$ is a separable, commutative, unital AF algebra. Its state space is the
weak-$\ast$ compact convex set $\Prob(\dis)$ of Borel probability measures on $\dis$,
and its condensation satisfies $\cond{C(\dis)}(S)=C(S\times\dis)$ for profinite $S$;
in particular $\dis$ itself is an admissible probe.
\end{proposition}

\begin{proof}
Each $\dis_{i}=Q^{\Lambda_{i}}$ is finite, so $C(\dis_{i})\cong\CC^{\abs{\dis_{i}}}$ is
finite-dimensional. The pullback maps $C(\dis_{i})\to C(\dis_{j})$ for $i\le j$ are
unital $\ast$-monomorphisms, and any $g\in C(\dis)$ is a uniform limit of functions
factoring through some $\dis_{i}$: the cylinder functions are dense by
Stone--Weierstrass, since they separate points of $\dis$ and are closed under the
algebra operations and conjugation. Hence the natural map
$\varinjlim_{i}C(\dis_{i})\to C(\dis)$ is an isometric $\ast$-isomorphism onto a
dense, hence closed and therefore full, subalgebra. This exhibits $C(\dis)$ as an AF
algebra; commutativity and unitality are clear, and separability follows as in
\Cref{prop:uhf}. By Riesz representation, states on $C(\dis)$ are Borel probability
measures. For the condensation, $\cond{C(\dis)}(S)=C(S,C(\dis))=C(S\times\dis)$ by the
exponential law for continuous maps out of the compact space $S$ into $C(\dis)$
(equivalently, the tensor identity $C(S)\otimes C(\dis)\cong C(S\times\dis)$ for
compact Hausdorff spaces).
\end{proof}

\subsection{The Bellissard crossed product and its condensation}

The observables of a homogeneous disordered system are not functions on $\dis$ alone
but the covariant operators; the algebra that carries the physics is the crossed
product.

\begin{progthm}[Crossed-product disorder algebra]
\label{thm:crossed}
Let $\dis=Q^{\ZZ^{d}}$ with the shift action $T$ of $\ZZ^{d}$. The covariant
observable algebra of a homogeneous disordered family is the crossed product
\[
  \dalg \;=\; C(\dis)\rtimes_{T}\ZZ^{d},
\]
a separable unital $C^{*}$-algebra. Its $K$-theory $K_{\ast}(\dalg)$ carries the
disorder-averaged topological invariants of the family, and an invariant ergodic
measure induces a trace whose range on $K_{0}$ is the gap-labelling group. The
condensation $\cond{\dalg}$ is functorial in the finite quotients of $\dis$: a
$T$-equivariant quotient $\dis\to\dis'$ induces a $\ast$-homomorphism
$\dalg'\to\dalg$ and hence a morphism $\cond{\dalg'}\to\cond{\dalg}$ of condensed
$C^{*}$-algebras, compatibly with composition.
\end{progthm}

\begin{proof}
$C(\dis)$ is a separable unital $C^{*}$-algebra (\Cref{prop:hullAF}) and $\ZZ^{d}$ is
countable discrete and \emph{amenable}, so the full and reduced crossed products
coincide and $\dalg=C(\dis)\rtimes_{T}\ZZ^{d}$ is an unambiguous separable unital
$C^{*}$-algebra by the standard construction; we write $\rtimes$ for this common
value. That its
$K$-theory computes the disorder-averaged invariants, and that a $T$-invariant
ergodic probability measure $\mathbb{P}$ on $\dis$ induces a faithful trace
$\tau_{\mathbb{P}}$ on $\dalg$ whose composition with the $K_{0}$-pairing is the
integrated density of states and yields the gap-labelling group, is Bellissard's
noncommutative geometry of the quantum Hall
effect~\mbox{\cite{bellissard-ncg-qhe,kellendonk-tilings,prodan-sb}}. For functoriality, a
$T$-equivariant continuous surjection $p:\dis\to\dis'$ induces a unital equivariant
$\ast$-monomorphism $p^{*}:C(\dis')\hookrightarrow C(\dis)$; by the universal
property of the crossed product an equivariant $\ast$-homomorphism of coefficient
algebras induces a $\ast$-homomorphism
$p^{*}\rtimes\ZZ^{d}:\dalg'=C(\dis')\rtimes_{T}\ZZ^{d}\to C(\dis)\rtimes_{T}\ZZ^{d}=\dalg$,
functorial in $p$. Applying condensation, which is a functor
(\Cref{prop:condcstar}), gives $\cond{\dalg'}\to\cond{\dalg}$, compatible with
composition. Since $\dis=\varprojlim_{i}\dis_{i}$, these assemble the condensation of
$\dalg$ over the inverse system of finite quotients.
\end{proof}

\begin{remark}[Bulk--boundary and coarse geometry]
The $K$-theory of $\dalg$ is not only a receptacle for a number; the bulk--boundary
correspondence realizes the boundary invariant as the image of the bulk class under
the connecting map of the Toeplitz extension of
$\dalg$~\cite{bourne-kellendonk-rennie}, and the coarse/Roe-algebraic picture gives
the same invariants a controlled-topology
description~\cite{kubota-controlled}. Each of these is a $K$-theory computation to
which \Cref{cor:condinvariant} applies, so the boundary invariant is condensed for
the same reason the bulk one is. The naturality of the \emph{solid} invariant in the
disorder hull (that solidification commutes with the maps above) is the content
of \Cref{conj:solidnat}.
\end{remark}

\section{Conjectures}
\label{sec:conjectures}

The theorems above are the provable core. The program's ambition on the
positivity/$C^{*}$ side is larger, and we state the open part as numbered
conjectures with stable identifiers. Each is separated cleanly from what is proved,
and each names the nearest rigorous result it would extend.

\begin{conjecture}[Positivity cuts out a closed condensed substack]
\label{conj:substack}
The assignment $H\mapsto(\qla,\St(\qla))$ upgrades to a morphism of condensed stacks
from $\Ham_{d,G}$ to the stack of condensed unital $C^{*}$-algebras equipped with
their condensed state space. Moreover, inside a larger ``formal-interaction'' stack
in which the $C^{*}$-relations are dropped, positivity together with the
$C^{*}$-identity $\norm{a^{*}a}=\norm{a}^{2}$ cuts out $\Ham_{d,G}$ as a
\emph{closed} condensed substack.
\end{conjecture}

\noindent This would promote \Cref{thm:statespace} from a statement about a single
algebra to one about the whole moduli object; the closedness is the stacky form of
``positivity is a closed condition,'' and its natural proof would run through the
condensed-stack technology of~\cite{clausen-scholze-analytic-stacks} together with
the condensed-set morphism $\RR\times\cond{\BF}\to\cond{\Aut(\qla)}$ of
Part~I~\cite{paperLocality}.

\begin{conjecture}[Naturality of the solid invariant in disorder]
\label{conj:solidnat}
For a disordered family $H\in\Ham_{d,G}(\dis)$ with observable algebra
$\dalg=C(\dis)\rtimes_{T}\ZZ^{d}$, the solid invariant
$\Solid(\Kalg(\cond{\dalg}))$ agrees with the operator-$K$ invariant $\Kop(\dalg)$
\emph{naturally in $\dis$}: the identifications of \Cref{thm:aoki} commute with the
$\ast$-homomorphisms induced by $T$-equivariant finite quotients of \Cref{thm:crossed},
so the invariant is compatible with every finite-resolution approximation of the
hull.
\end{conjecture}

\noindent This is the disorder-averaged, family version of \Cref{cor:condinvariant}.
It is the precise sense in which ``profinite probes match the configuration space'':
the invariant is not merely definable over $\dis$ but is the limit of its values on
the finite quotients $\dis_{i}$.

\Needspace{9\baselineskip}
\begin{conjecture}[Real/$KKO$ refinement of the periodic table]
\label{conj:real}
The $KO$-graded Kitaev periodic table is recovered by solidification of the
\emph{Real} (in the sense of Atiyah's $KR$) algebraic $K$-theory of the condensed
real observable algebra: for each of the ten symmetry classes, the class-appropriate
$KR$-degree of $\Solid(\Kalg(\cond{A}))$ equals the operator $KKO$-invariant of the
system, and the eightfold real periodicity matches the eightfold structure of the
real classes.
\end{conjecture}

\noindent \Cref{thm:aoki} is a real, not yet a Real, statement, and the tenfold way
needs the $KR$ grading that keeps track of the anti-linear time-reversal and
particle--hole symmetries~\cite{prodan-sb,bourne-kellendonk-rennie,thiang-ktheory}.
The conjecture is that Aoki's solidification is compatible with the Real structure so
that the whole table, not just its complex and real-linear parts, is condensed.

\begin{conjecture}[Solid state space as a classifying object]
\label{conj:classifying}
The condensed state space $\cond{\St(\qla)}$, with the convex structure making it a
solid $\RR$-convex object, is a classifying object for condensed cyclic
representations: a continuous family of states over a profinite base $S$ is the same
as a map $S\to\cond{\St(\qla)}$, and the GNS functor of \Cref{prop:gns} descends to a
universal condensed field of representations over $\cond{\St(\qla)}$, refining the
continuity of \Cref{thm:groundstate} to a continuous field of Hilbert spaces on the
locus where the GNS dimension is locally constant.
\end{conjecture}

\noindent This is the condensed form of the state-space approach to phase
classification, in which phases are homotopy classes of sections of bundles of
(pure) states~\cite{denittis-states,denittis-rendel-weyl}; the conjecture is that
those bundles and their sections are the shadow of a single condensed classifying
object, and that GNS is its universal representation.

\begin{conjecture}[Condensed enhancement of the split property and superselection]
\label{conj:split}
The split property of gapped ground states and the
Doplicher--Haag--Roberts superselection structure admit condensed enhancements: the
sector structure --- a symmetric or braided tensor category in the relevant
dimension --- is detected by a condensed invariant over profinite bases, and the
split property provides the condensed factorization along which $\pi_{0}$ of the
phase stack descends through finite quotients of a disorder hull.
\end{conjecture}

\noindent The split property is central to the operator-algebraic classification of
gapped phases and to the construction of the indices that make that classification
work~\cite{ogata-classification-review,ogata-h3-index}. The conjecture asks that this
structure be visible to the profinite probes, so that the descent of phase labels
along disorder approximations (the theme of \Cref{thm:crossed} and
\Cref{conj:solidnat}) has a representation-theoretic counterpart.

\section{Discussion and limitations}
\label{sec:discussion}

\paragraph{Scope of the solid bridge.}
\Cref{thm:aoki} and \Cref{cor:condinvariant} reach the invariants of free-fermion
and disordered phases (the invertible, $K$-theoretic sector). They say nothing
about noninvertible topological order, whose excitations form a braided or modular
tensor category rather than a $K$-theory class. A condensed account of that order
would need a condensed \emph{higher} stack of phases and defects, not a $K$-theory
spectrum, and it is out of scope here; Parts~V and~VI record the boundary of what the
$K$-theoretic story can classify. Within its scope, the bridge is exact, but its
scope is the invertible sector.

\paragraph{Real versus complex, connective versus periodic.}
We have been deliberate about Aoki's hypotheses. The theorem is about real algebras;
the operator $K$-theory it recovers is $KO$-flavored, which is the right flavor for
the periodic table, but the identifications carry connective-versus-periodic and
real-versus-Real bookkeeping that we have imported rather than restated. A reader who
wants the graded-group form must read it off the source with those conventions in
place. Overstating the theorem as a naked Bott-periodic isomorphism would be a
mistake, and we have avoided it.

\paragraph{Unital versus nonunital.}
Our systems have finite on-site dimension, so $\qla$ is unital and $\St(\qla)$ is
compact: the hypotheses of \Cref{lem:alaoglu} and \Cref{thm:statespace} hold on
the nose. For systems whose observable algebra is nonunital (for instance certain
half-space or continuum models) the state space is only weak-$\ast$ locally compact,
and one must pass to the multiplier algebra or unitization before condensing;
\Cref{thm:statespace} then applies to the unitization. We flag this rather than hide
it.

\paragraph{Light versus full condensed.}
Separability of $\qla$ and of $C(\dis)$ lets us work with light condensed sets
throughout, which sidesteps the set-theoretic size issues of the unrestricted theory
and is exactly matched to the countable inverse limits that describe finite-resolution
disorder data. For systems that are genuinely non-separable this convenience is lost
and the full condensed formalism, with its universe bookkeeping, is required.

\paragraph{What is genuinely new here.}
Nothing in \Cref{sec:framework,sec:states,sec:gns,sec:disorder} is a new theorem of
operator algebras; the individual facts are classical or cited. The new element is
their assembly: the state space, positivity, GNS, and the crossed product are placed
in the condensed category so that Aoki's bridge can act on them, and the ground-state
section is shown to be a genuine condensed point over the gapped locus. The synthesis
is the contribution, and its honest measure is that it makes an existing invariant
condensed, not that it makes a new one.

\section{Conclusion}
\label{sec:conclusion}

We have carried the observable side of a quantum lattice system (its quasi-local
$C^{*}$-algebra, its positivity structure, its state space, and its $K$-theory)
into condensed mathematics, and we have been explicit about the line between theorem
and conjecture. On the theorem side: the quasi-local algebra is a separable AF
$C^{*}$-algebra whose condensation is a light sheaf of $C^{*}$-algebras
(\Cref{prop:uhf,prop:condcstar}); the state space condenses to a compact Hausdorff
condensed set on which positivity is a closed condition (\Cref{thm:statespace}); GNS
is a functor and the ground-state section over a uniformly gapped family is
weak-$\ast$ continuous (\Cref{prop:gns,thm:groundstate}); Aoki's solidification
theorem makes the operator-$K$ invariant of a real Banach observable algebra a
condensed-mathematical object (\Cref{thm:aoki,cor:condinvariant}); and the disorder
crossed product condenses functorially in the finite quotients of the profinite hull
(\Cref{prop:hullAF,thm:crossed}). On the conjecture side we have named five open
problems (\Cref{conj:substack,conj:solidnat,conj:real,conj:classifying,conj:split})
whose resolution would complete the positivity/$C^{*}$ module of the program.

The single load-bearing insight is worth isolating. A topological invariant of a
free-fermion or disordered phase is an operator-$K$ class, and by Aoki's theorem that
class is the Bott-inverted solidification of the connective algebraic $K$-theory of a
condensed algebra. The invariant is unchanged; its provenance is now condensed, and
with that provenance comes the family-theoretic and profinite-disorder machinery that
the rest of the program deploys. That is the whole of what this paper claims, and it
is stated in the key of a research program rather than a finished theory.

\section*{Code availability}
The formal-verification code accompanying this paper (finite-dimensional
$C^{*}$-algebra positivity checks, the explicit GNS construction, and the Bloch-ball
geometry of qubit states) is at
\href{https://github.com/YonedaAI/topological-phases-of-matter}{github.com/YonedaAI/topological-phases-of-matter},
in \texttt{src/positivity-cstar-norms/}.

\begin{thebibliography}{99}

\bibitem{clausen-scholze-condensed}
P.~Scholze (joint with D.~Clausen),
\emph{Lectures on Condensed Mathematics},
stable citable version of the 2019 Bonn course, \href{https://arxiv.org/abs/2605.03658}{arXiv:2605.03658}.

\bibitem{clausen-scholze-analytic}
P.~Scholze (joint with D.~Clausen),
\emph{Lectures on Analytic Geometry},
stable citable version of the 2019/20 Bonn course, \href{https://arxiv.org/abs/2605.03655}{arXiv:2605.03655}.

\bibitem{clausen-scholze-analytic-stacks}
D.~Clausen and P.~Scholze,
\emph{Analytic Stacks},
lecture course, Max Planck Institute for Mathematics / IHES, 2023--2024.

\bibitem{barwick-haine-pyknotic}
C.~Barwick and P.~Haine,
\emph{Pyknotic objects, I. Basic notions},
\href{https://arxiv.org/abs/1904.09966}{arXiv:1904.09966} (2019).

\bibitem{aoki-solidification}
K.~Aoki,
\emph{(Semi)topological $K$-theory via solidification},
\href{https://arxiv.org/abs/2409.01462}{arXiv:2409.01462} (2024).

\bibitem{bmns-automorphic}
S.~Bachmann, S.~Michalakis, B.~Nachtergaele, and R.~Sims,
\emph{Automorphic equivalence within gapped phases of quantum lattice systems},
Comm. Math. Phys. \textbf{309} (2012), 835--871, \href{https://arxiv.org/abs/1102.0842}{arXiv:1102.0842}.

\bibitem{hastings-wen-qac}
M.~B.~Hastings and X.-G.~Wen,
\emph{Quasi-adiabatic continuation of quantum states: the stability of topological
ground state degeneracy and emergent gauge invariance},
Phys. Rev. B \textbf{72} (2005), 045141, \href{https://arxiv.org/abs/cond-mat/0503554}{arXiv:cond-mat/0503554}.

\bibitem{nsy-quasilocality-1}
B.~Nachtergaele, R.~Sims, and A.~Young,
\emph{Quasi-locality bounds for quantum lattice systems. I. Lieb--Robinson bounds,
quasi-local maps, and spectral flow automorphisms},
J. Math. Phys. \textbf{60} (2019), 061101, \href{https://arxiv.org/abs/1810.02428}{arXiv:1810.02428}.

\bibitem{nsy-bulkgap}
B.~Nachtergaele, R.~Sims, and A.~Young,
\emph{Stability of the bulk gap for frustration-free topologically ordered quantum
lattice systems},
Lett. Math. Phys. \textbf{114} (2024), 24, \href{https://arxiv.org/abs/2102.07209}{arXiv:2102.07209}.

\bibitem{kitaev-periodic-table}
A.~Kitaev,
\emph{Periodic table for topological insulators and superconductors},
AIP Conf. Proc. \textbf{1134} (2009), 22--30, \href{https://arxiv.org/abs/0901.2686}{arXiv:0901.2686}.

\bibitem{thiang-ktheory}
G.~C.~Thiang,
\emph{On the $K$-theoretic classification of topological phases of matter},
Ann. Henri Poincar\'e \textbf{17} (2016), 757--794, \href{https://arxiv.org/abs/1406.7366}{arXiv:1406.7366}.

\bibitem{prodan-sb}
E.~Prodan and H.~Schulz-Baldes,
\emph{Bulk and Boundary Invariants for Complex Topological Insulators: From
$K$-Theory to Physics},
Mathematical Physics Studies, Springer (2016), \href{https://arxiv.org/abs/1510.08744}{arXiv:1510.08744}.

\bibitem{ssh-1979}
W.~P.~Su, J.~R.~Schrieffer, and A.~J.~Heeger,
\emph{Solitons in polyacetylene},
Phys. Rev. Lett. \textbf{42} (1979), 1698--1701.

\bibitem{tknn-1982}
D.~J.~Thouless, M.~Kohmoto, M.~P.~Nightingale, and M.~den Nijs,
\emph{Quantized Hall conductance in a two-dimensional periodic potential},
Phys. Rev. Lett. \textbf{49} (1982), 405--408.

\bibitem{bellissard-ncg-qhe}
J.~Bellissard, A.~van Elst, and H.~Schulz-Baldes,
\emph{The noncommutative geometry of the quantum Hall effect},
J. Math. Phys. \textbf{35} (1994), 5373--5451, \href{https://arxiv.org/abs/cond-mat/9411052}{arXiv:cond-mat/9411052}.

\bibitem{kellendonk-tilings}
J.~Kellendonk,
\emph{Non commutative geometry of tilings and gap labelling},
Rev. Math. Phys. \textbf{7} (1995), 1133--1180, \href{https://arxiv.org/abs/cond-mat/9403065}{arXiv:cond-mat/9403065}.

\bibitem{bourne-kellendonk-rennie}
C.~Bourne, J.~Kellendonk, and A.~Rennie,
\emph{The $K$-theoretic bulk-edge correspondence for topological insulators},
Ann. Henri Poincar\'e \textbf{18} (2017), 1833--1866, \href{https://arxiv.org/abs/1604.02337}{arXiv:1604.02337}.

\bibitem{kubota-controlled}
Y.~Kubota,
\emph{Controlled topological phases and bulk-edge correspondence},
Comm. Math. Phys. \textbf{349} (2017), 493--525, \href{https://arxiv.org/abs/1511.05314}{arXiv:1511.05314}.

\bibitem{ogata-classification-review}
Y.~Ogata,
\emph{Classification of gapped ground state phases in quantum spin systems},
Proc. ICM 2022, \href{https://arxiv.org/abs/2110.04675}{arXiv:2110.04675}.

\bibitem{ogata-h3-index}
Y.~Ogata,
\emph{A $H^{3}(G,\mathbb{T})$-valued index of symmetry protected topological phases
with on-site finite group symmetry for two-dimensional quantum spin systems},
Forum Math. Pi \textbf{9} (2021), e13, \href{https://arxiv.org/abs/2101.00426}{arXiv:2101.00426}.

\bibitem{denittis-states}
G.~De Nittis,
\emph{Topological phases of non-interacting systems: a general approach based on
states},
Lett. Math. Phys. \textbf{115} (2025), \href{https://arxiv.org/abs/2502.03590}{arXiv:2502.03590}.

\bibitem{denittis-rendel-weyl}
G.~De Nittis and S.~G.~Rendel,
\emph{A scheme for topological phases of the Weyl $C^{*}$-algebra},
\href{https://arxiv.org/abs/2607.01183}{arXiv:2607.01183} (2026).

\bibitem{paperLocality}
M.~Long and The YonedaAI Collaboration,
\emph{Condensed Locality: Lieb--Robinson Estimates and Quasi-Local Dynamics on the
Moduli Stack of Hamiltonians},
Part~I of this series (2026).

\bibitem{paperGap}
M.~Long and The YonedaAI Collaboration,
\emph{The Uniformly Gapped Substack: Existence and Stability of the Thermodynamic
Spectral Gap},
Part~III of this series (2026).

\bibitem{paperEFT}
M.~Long and The YonedaAI Collaboration,
\emph{From Lattice Models to Effective Field Theories: Stabilization and the
Invertible Condensed Phase Spectrum},
Part~IV of this series (2026).

\bibitem{paperRealizability}
M.~Long and The YonedaAI Collaboration,
\emph{Physical Realizability of Bordism and Homotopy Classes by Gapped Lattice
Systems},
Part~V of this series (2026).

\bibitem{paperSynthesis}
M.~Long and The YonedaAI Collaboration,
\emph{Topological Phases of Matter in the Condensed-Mathematics Paradigm: A Modular
Research Program},
Part~VI of this series (2026).

\end{thebibliography}

\end{document}
